\documentclass[11pt]{article}
\usepackage[utf8]{inputenc}
\usepackage{amsmath,amssymb}
\usepackage{hyperref}
\title{Efficient Brain Computer Interface Decoding: A Resource-Aware Approach}
\author{Anonymous}
\date{}
\begin{document}
\maketitle

\begin{abstract}
This paper studies Brain Computer Interface Decoding. Grounded in a real, citable literature \cite{ref1,ref2,ref3,ref4,ref5,ref6,ref7,ref8},
we propose a focused method and report \textbf{illustrative} experiments.
All references below are real and verifiable (OpenAlex/arXiv); the reported
numbers are simulated to illustrate intended behaviour.
\end{abstract}

\section{Introduction}
Brain Computer Interface Decoding is an active research area. This work targets a concrete gap identified from
the recent literature and contributes a method together with an evaluation protocol.

\section{Related Work (real literature)}
The following works, retrieved from public bibliographic APIs, frame our study:
\begin{itemize}
\item Hochberg et al. (2006) — "Neuronal ensemble control of prosthetic devices by a human with tetraplegia", Nature (cited 3338).
\item Hochberg et al. (2012) — "Reach and grasp by people with tetraplegia using a neurally controlled robotic arm", Nature (cited 2750).
\item Epanechnikov (1969) — "Non-Parametric Estimation of a Multivariate Probability Density", Theory of Probability and Its Applications (cited 1838).
\item Willett et al. (2021) — "High-performance brain-to-text communication via handwriting", Nature (cited 968).
\item Abiri et al. (2018) — "A comprehensive review of EEG-based brain–computer interface paradigms", Journal of Neural Engineering (cited 955).
\item O’Sullivan et al. (2014) — "Attentional Selection in a Cocktail Party Environment Can Be Decoded from Single-Trial EEG", Cerebral Cortex (cited 924).
\end{itemize}

\section{Method}
We reduce the dominant computational cost via adaptive resource allocation, activating only the components needed per input. We instantiate this idea for Brain Computer Interface Decoding and analyse its cost/quality trade-off.

\section{Experiments (Illustrative)}
\begin{center}
\begin{tabular}{lcc}
System & Primary Metric & Efficiency \\ \hline
Strong baseline & 0.812 & 1.00$\times$ \\
Ablation & 0.828 & 0.92$\times$ \\
\textbf{Ours} & \textbf{0.851} & \textbf{0.71$\times$}
\end{tabular}
\end{center}
\emph{Metrics simulated to illustrate the intended effect; not measured.}

\section{Conclusion}
We connected a real literature to a concrete method for Brain Computer Interface Decoding. Future work will
validate the illustrative results with real experiments.

\begin{thebibliography}{9}
\bibitem{ref1} Leigh R. Hochberg, Mijail D. Serruya, Gerhard M. Friehs et al.. Neuronal ensemble control of prosthetic devices by a human with tetraplegia. \emph{Nature}, 2006. DOI:10.1038/nature04970
\bibitem{ref2} Leigh R. Hochberg, Daniel Bacher, Beata Jarosiewicz et al.. Reach and grasp by people with tetraplegia using a neurally controlled robotic arm. \emph{Nature}, 2012. DOI:10.1038/nature11076
\bibitem{ref3} V. A. Epanechnikov. Non-Parametric Estimation of a Multivariate Probability Density. \emph{Theory of Probability and Its Applications}, 1969. DOI:10.1137/1114019
\bibitem{ref4} Francis R. Willett, Donald T. Avansino, Leigh R. Hochberg et al.. High-performance brain-to-text communication via handwriting. \emph{Nature}, 2021. DOI:10.1038/s41586-021-03506-2
\bibitem{ref5} Reza Abiri, Soheil Borhani, Eric W. Sellers et al.. A comprehensive review of EEG-based brain–computer interface paradigms. \emph{Journal of Neural Engineering}, 2018. DOI:10.1088/1741-2552/aaf12e
\bibitem{ref6} James O’Sullivan, Alan J. Power, Nima Mesgarani et al.. Attentional Selection in a Cocktail Party Environment Can Be Decoded from Single-Trial EEG. \emph{Cerebral Cortex}, 2014. DOI:10.1093/cercor/bht355
\bibitem{ref7} Thorsten O. Zander, Christian Kothe. Towards passive brain–computer interfaces: applying brain–computer interface technology to human–machine systems in general. \emph{Journal of Neural Engineering}, 2011. DOI:10.1088/1741-2560/8/2/025005
\bibitem{ref8} Tijl Grootswagers, Susan G. Wardle, Thomas A. Carlson. Decoding Dynamic Brain Patterns from Evoked Responses: A Tutorial on Multivariate Pattern Analysis Applied to Time Series Neuroimaging Data. \emph{Journal of Cognitive Neuroscience}, 2016. DOI:10.1162/jocn\_a\_01068
\end{thebibliography}
\end{document}
